% CV de Valentino Mensio. Fuente del PDF publicado en public/CV-Valentino-Mensio.pdf.
% Se compila en Overleaf (pdfLaTeX) junto a la carpeta cv-icons/ (mail, phone, linkedin,
% github, web .png). Datos y cifras alineados con src/data/profile.ts (2026-11-25).
\documentclass[letterpaper,10pt]{article}
\usepackage[margin=0.38in]{geometry}
\usepackage[spanish]{babel}
\usepackage[utf8]{inputenc}
\usepackage[T1]{fontenc}
\usepackage{enumitem}
\usepackage[hidelinks]{hyperref}
\usepackage{titlesec}
\usepackage{xcolor}
\usepackage{tabularx}
\usepackage{array}
\usepackage{ragged2e}
\usepackage{microtype}
\usepackage{graphicx}
\linespread{1.0}

\definecolor{accent}{HTML}{1F4E79}
\definecolor{textgray}{HTML}{222222}
\hypersetup{colorlinks=true, urlcolor=accent, linkcolor=accent,
  pdftitle={Valentino Mensio - Desarrollador Backend / Full Stack},
  pdfauthor={Valentino Mensio},
  pdfkeywords={Python, FastAPI, PostgreSQL, React, TypeScript, Docker, Backend, Full Stack}}

\pagestyle{empty}
\renewcommand{\familydefault}{\sfdefault}
\setlength{\parindent}{0pt}
\setlength{\parskip}{0.8pt}
\color{textgray}

\titleformat{\section}
  {\large\bfseries\scshape\raggedright\color{accent}}
  {}{0em}{}[\vspace{2pt}{\color{accent}\titlerule[0.8pt]}\vspace{-3pt}]
\titlespacing*{\section}{0pt}{4pt}{2pt}

\newcommand{\entry}[4]{%
  \textbf{#1}\hfill\textit{#2}\\[-1pt]
  \textit{#3}\if\relax\detokenize{#4}\relax\else\hfill{\footnotesize\textbf{Stack:} #4}\fi\par\vspace{0.5pt}
}

% Íconos: imágenes PNG, no glifos de fuente, para que los lectores automáticos (ATS)
% no extraigan caracteres basura. Los enlaces muestran la dirección, no sólo el nombre.
\newcommand{\contacticon}[1]{\raisebox{-0.18em}{\includegraphics[height=0.95em]{cv-icons/#1.png}}}
\newcommand{\contactitem}[3]{\contacticon{#1}\,\href{#2}{#3}}
\newcommand{\bloglink}[1]{{\footnotesize\contacticon{web}\,\href{#1}{Caso técnico}}}

\setlist[itemize]{leftmargin=1.05em, itemsep=0.75pt, topsep=1.5pt, parsep=0pt, partopsep=0pt}

\begin{document}

\begin{center}
  {\Huge\bfseries\scshape Valentino Mensio}\\[-1pt]
  {\large Desarrollador Backend / Full Stack}\\[3pt]
  \small
  \contactitem{mail}{mailto:valentinomensio@gmail.com}{valentinomensio@gmail.com} \quad | \quad
  \contactitem{phone}{tel:+5493547608293}{+54 9 3547 608293} \quad | \quad
  Córdoba, Argentina · remoto o híbrido\\[2pt]
  \contactitem{linkedin}{https://www.linkedin.com/in/valentino-mensio}{linkedin.com/in/valentino-mensio} \quad | \quad
  \contactitem{github}{https://github.com/ValentinoMensio}{github.com/ValentinoMensio} \quad | \quad
  \contactitem{web}{https://valentinomensio.com}{valentinomensio.com}
\end{center}
\vspace{-7pt}

%-----------PERFIL PROFESIONAL-----------
\section{Perfil profesional}
\vspace{2pt}
\begingroup\linespread{1.04}\selectfont
\textbf{Analista en Computación por FAMAF (UNC)} y desarrollador \textbf{backend / full stack}. Desarrollo y opero \textbf{dos sistemas propios en producción} que usan clientes reales todos los días: gestión y \textbf{facturación electrónica con ARCA} para una fábrica de alimentos, y un \textbf{producto modular de gestión comercial con punto de venta}. Me especializo en \textbf{Python, FastAPI, PostgreSQL}, diseño de APIs y modelado de dominio, con foco en consistencia de datos, \textbf{testing} y despliegue con \textbf{Docker}. Busco sumarme a un equipo donde aportar autonomía y criterio técnico, y seguir aprendiendo.
\endgroup

%-----------EXPERIENCIA-----------
\section{Experiencia}
\vspace{2pt}
\entry{Gestión · producto propio \bloglink{https://valentinomensio.com/blog/gestion-arquitectura-modular/}}{Julio 2026 -- actualidad}{Desarrollador Backend / Full Stack}{Python, FastAPI, PostgreSQL, React, TypeScript, Docker}
\begin{itemize}
  \item Llevé a \textbf{producción en 4 meses} (14 versiones) un producto de gestión comercial con \textbf{19 módulos}: ventas, inventario, POS, facturación, cuentas corrientes, tesorería y consignación. Desde noviembre de 2026 opera en un \textbf{comercio minorista}.
  \item Diseñé un \textbf{monolito modular} con contratos públicos entre dominios, verificados por tests de arquitectura; API \textbf{OpenAPI} con cliente TypeScript generado y 177 migraciones Alembic.
  \item Resolví reintentos y fallos parciales con \textbf{idempotencia} persistente, transacciones acotadas, outbox con worker y un \textbf{checkout POS reanudable} que no duplica cobros.
  \item \textbf{Más de 3.000 tests} (unitarios, integración con PostgreSQL real, contratos y E2E con Playwright), CI y \textbf{más de 190 pull requests}.
  \item Despliegue por cliente con imágenes Docker versionadas, perfil de módulos, backup previo a cada actualización y Cloudflare Tunnel.
\end{itemize}

\entry{Granalia · cliente \bloglink{https://valentinomensio.com/blog/granalia-sistema-facturacion/}}{Abril 2026 -- actualidad}{Desarrollador Full Stack}{FastAPI, React, PostgreSQL, Docker, ARCA}
\begin{itemize}
  \item Desarrollé el sistema de gestión comercial que una \textbf{fábrica de alimentos usa a diario desde abril de 2026}: reemplazó planillas para pedidos, clientes, listas de precios, remitos, compras y stock.
  \item Integré \textbf{facturación electrónica con ARCA} (WSAA / WSFEv1): Factura A y notas de crédito con CAE, recuperación ante timeouts \textbf{sin duplicar numeración}, comprobantes autorizados bloqueados y auditoría sin secretos.
  \item Construí estadísticas para la dirección: \textbf{rentabilidad por producto} con costos de materia prima y envases, costo de reposición, rankings y mapa de clientes.
  \item Me ocupo del desarrollo, el despliegue en \textbf{Ubuntu Server} (Docker, PostgreSQL, Caddy) y el soporte continuo.
\end{itemize}

%-----------PROYECTOS-----------
\section{Proyectos}
\vspace{2pt}
\entry{BeLeadAI · producto propio \bloglink{https://valentinomensio.com/blog/beleadai-automatizacion-prospeccion/}}{Noviembre 2025 -- marzo 2026}{Desarrollador Full Stack}{FastAPI, MySQL, Redis, Selenium, Chrome MV3}
\begin{itemize}
  \item Plataforma de \textbf{prospección automatizada} en Instagram: una extensión Chrome conectada a una API privada extrae audiencias, analiza perfiles y encola campañas.
  \item Diseñé \textbf{jobs asíncronos}, colas, workers por cuenta, cuotas, reintentos, idempotencia y seguimiento de estados por WebSockets, contemplando errores de scraping y sesiones.
\end{itemize}

\entry{El Switcher · FAMAF, equipo \bloglink{https://valentinomensio.com/blog/el-switcher-backend-scrum/}}{Septiembre -- octubre 2024}{Desarrollador Backend y Scrum Master}{FastAPI, WebSockets, React, Pytest}
\begin{itemize}
  \item Implementé reglas de juego y partidas \textbf{multijugador en tiempo real} con WebSockets y \textbf{arquitectura hexagonal}; el backend alcanzó \textbf{95\% de cobertura} de tests.
  \item Fui \textbf{Scrum Master} en el Sprint 2: planificación, dailies, tickets con criterios de aceptación y seguimiento de bloqueos.
\end{itemize}

%-----------EDUCACIÓN-----------
\section{Educación}
\vspace{2pt}
\entry{Universidad Nacional de Córdoba -- FAMAF}{2021 -- actualidad}{Analista en Computación (título obtenido) · Licenciatura en Ciencias de la Computación en curso}{}
\entry{Instituto Técnico Superior ``El Obraje''}{2014 -- 2020}{Técnico en Informática Profesional y Personal}{}

%-----------HABILIDADES-----------
\section{Habilidades técnicas}
\vspace{2pt}
\begin{tabularx}{\textwidth}{@{}>{\bfseries}p{0.22\textwidth}X@{}}
Lenguajes: & Python | TypeScript | JavaScript | SQL \\
Backend y APIs: & FastAPI | Pydantic | SQLAlchemy | REST | OpenAPI | WebSockets | Flask \\
Datos: & PostgreSQL | Alembic | MySQL | Redis | SQLite | MinIO / S3 \\
Frontend: & React | TypeScript | TanStack Query | Vite | Tailwind CSS | Chrome Extensions MV3 \\
Calidad: & Pytest | Playwright | Vitest | mypy | Ruff | tests de arquitectura y contratos \\
DevOps: & Docker / Compose | Linux / Ubuntu Server | GitHub Actions | Caddy | Nginx | Cloudflare Tunnel | backups y restore \\
Integraciones: & ARCA (WSAA / WSFEv1) | OpenAI API | Selenium | generación de PDF \\
Proceso: & Git | GitHub (issues, PRs, releases) | Jira | Scrum \\
Idiomas: & Español nativo | Inglés básico (en formación) \\
\end{tabularx}

\end{document}
